\documentclass{article}

\usepackage{geometry}
\usepackage{makecell}
\usepackage{array}
\usepackage{multicol}
\usepackage{setspace}
\usepackage{changepage}
\usepackage{booktabs}
\usepackage[explicit]{titlesec}
\usepackage{hyperref}
\usepackage{graphicx}
\usepackage{cprotect}
\usepackage{float}
\newcolumntype{?}{!{\vrule width 1pt}}
\newcommand{\paragraphlb}[1]{\paragraph{#1}\mbox{}\\}
\renewcommand{\contentsname}{Inhaltsverzeichnis:}
\renewcommand\theadalign{tl}
\setstretch{1.10}
\setlength{\parindent}{0pt}

\titleformat{\section}
  {\normalfont\Large\bfseries}{\thesection}{1em}{\hyperlink{sec-\thesection}{#1}
\addtocontents{toc}{\protect\hypertarget{sec-\thesection}{}}}
\titleformat{name=\section,numberless}
  {\normalfont\Large\bfseries}{}{0pt}{#1}

\titleformat{\subsection}
  {\normalfont\large\bfseries}{\thesubsection}{1em}{\hyperlink{subsec-\thesubsection}{#1}
\addtocontents{toc}{\protect\hypertarget{subsec-\thesubsection}{}}}
\titleformat{name=\subsection,numberless}
  {\normalfont\large\bfseries}{\thesubsection}{0pt}{#1}

\hypersetup{
    colorlinks,
    citecolor=black,
    filecolor=black,
    linkcolor=black,
    urlcolor=black
}

\geometry{top=12mm, left=1cm, right=2cm}
\title{\vspace{-1cm}Digitale Geschäftsmodelle 1 - Netzeffekte - Apple}
\author{Andreas Hofer, Patrick Pramberger}

\begin{document}
	\maketitle
	\begin{enumerate}
		\item{App Store}
		\begin{itemize}
			\item{Der App Store hat einen indirekten Netzeffekt, da die Anzahl der Benutzer diesen für Entwickler attraktiver macht, was ihn wiederum für Benutzer attraktiver macht.}
		\end{itemize}
		\item{AirDrop, iMessage, FaceTime, iCloud}
		\begin{itemize}
			\item{Da man all diese Dienste nur verwenden kann, wenn beide Parteien ein Produkt von Apple besitzt, nimmt der Wert davon mit der Anzahl an Benutzern zu.}
		\end{itemize}
		\item{ApplePay}
		\begin{itemize}
			\item{ApplePay kann in vielen Webshops als Zahlungsmittel verwendet werden, wodurch ein indirekter Netzeffekt entsteht.}
		\end{itemize}
		\item{AirTag}
		\begin{itemize}
			\item{AirTags verwenden Applegeräte in der Umgebung um sich mit dem Server zu verbinden und so ihren Standort zu teilen, wodurch ein direkter Netzeffekt entsteht.}
		\end{itemize}
	\end{enumerate}
	























  
\end{document}